\documentclass[11pt]{amsart}
\usepackage[margin=1.2in]{geometry}
\usepackage{amsmath,amssymb,amsthm}
\usepackage[colorlinks,linkcolor=blue,citecolor=blue,urlcolor=blue]{hyperref}

\newtheorem{theorem}{Theorem}
\newtheorem{proposition}[theorem]{Proposition}
\newtheorem{lemma}[theorem]{Lemma}
\theoremstyle{remark}
\newtheorem{remark}[theorem]{Remark}
\newtheorem{question}[theorem]{Question}

\DeclareMathOperator{\per}{per}
\newcommand{\half}{\tfrac12}

\title[Sparsity is not enough]{Sparsity is not enough: Foregger's inequality fails\\
for fully indecomposable matrices}
\author{Tom Maciejewski}
\address{Independent researcher, New Jersey, USA}
\email{tommaciejewski4@gmail.com}
\subjclass[2020]{15A15, 15B51, 05C50}
\keywords{permanent, doubly stochastic matrix, fully indecomposable matrix,
nearly decomposable matrix, Foregger's conjecture}

\begin{document}

\begin{abstract}
A nearly decomposable matrix of order $n\ge3$ has at most $3n-3$ nonzero entries,
and if it is doubly stochastic then its permanent is at least $2^{1-n}$. This
is Foregger's inequality, proved by Lavi in 2026. Is it the sparsity of these
matrices that matters, or their minimality? We show that sparsity is not
enough. There is a fully indecomposable doubly stochastic $6\times6$ matrix
with $15=3n-3$ positive entries, all equal to $\frac14$, $\frac38$ or
$\frac12$, whose permanent is $\frac{15}{512}<\frac{16}{512}$; checking this
takes three lines, and no smaller order allows it. The threshold is exact: a
fully indecomposable doubly stochastic matrix with at most $2n+2$ positive
entries satisfies the inequality, and in every order $n\ge3$ the same
construction gives one with $2n+3$ positive entries that does not. Chaining the
construction gives matrices within the $3n-3$ budget whose permanents fall below
$2^{1-n}$ by a factor that is exponentially small in $n$.
\end{abstract}

\maketitle

\section{Introduction}

A $(0,1)$-matrix of order $n\ge2$ is \emph{fully indecomposable} if it has no
$p\times(n-p)$ zero submatrix with $1\le p\le n-1$, and \emph{nearly
decomposable} if it is fully indecomposable and stops being so whenever one of
its $1$s is replaced by $0$. A nonnegative matrix has either property when its
zero pattern does. Let $\Omega_n$ be the set of doubly stochastic matrices of
order $n$.

Foregger~\cite{Foregger} conjectured in 1980 that every nearly decomposable
$A\in\Omega_n$ satisfies
\begin{equation}\label{eq:foregger}
  \per A\ \ge\ 2^{1-n},
\end{equation}
with equality only for the matrices $P\cdot\half(I+C_n)\cdot Q$, where $C_n$
is the permutation matrix of an $n$-cycle and $P,Q$ are permutation matrices.
Lavi~\cite{Lavi} proved this in 2026.

By a theorem of Minc~\cite{Minc}, a nearly decomposable matrix of order
$n\ge3$ has at most $3n-3$ positive entries. So~\eqref{eq:foregger} is a
statement about very sparse doubly stochastic matrices, and it is natural to
ask whether the sparsity alone is responsible.
The following question was raised by T.~H. Foregger in conversation with the author.

\begin{quote}
\emph{Does every $A\in\Omega_n$ with at most $3n-3$ positive entries satisfy
$\per A\ge2^{1-n}$?}
\end{quote}

Without a bound on the number of entries the answer is of course no. The
$3\times3$ matrix with all entries $\frac13$ has permanent
$\frac29<\frac14$, and in general the van der Waerden minimum $n!/n^n$ is far
below $2^{1-n}$. But that matrix has $9$ positive entries where only
$3n-3=6$ are allowed, and it is not obvious that a matrix within the budget
can afford a dense block. It can, as soon as $n=6$.

\begin{theorem}\label{thm:six}
The matrix
\[
  A=\frac18\begin{pmatrix}
   3&3&2&0&0&0\\
   3&3&2&0&0&0\\
   2&2&0&4&0&0\\
   0&0&0&4&4&0\\
   0&0&0&0&4&4\\
   0&0&4&0&0&4
  \end{pmatrix}
\]
is doubly stochastic and fully indecomposable, has $15=3\cdot6-3$ positive
entries, and satisfies
\[
  \per A=\frac{15}{512}<\frac{16}{512}=2^{1-6}.
\]
\end{theorem}

\begin{proof}[Verification of the permanent]
Consider the permutations that contribute to $\per(8A)$. Those that use
position $(4,4)$ must also use $(5,5)$ and $(6,6)$. Together they contribute
$4^3\cdot48$, where $48$ is the permanent of the leading $3\times3$ block. The
others use $(4,5)$, hence $(5,6)$, $(6,3)$ and $(3,4)$. They contribute
$4^4\cdot18$, where $18$ is the permanent of the leading $2\times2$ block. So
$\per(8A)=3072+4608=7680$, whereas $8^6\cdot2^{-5}=8192$. The other properties
are proved in Lemma~\ref{lem:necklace}.
\end{proof}

The same idea works in every order $n\ge3$ with only $2n+3$ positive entries,
which is within the budget $3n-3$ exactly when $n\ge6$.

\begin{theorem}\label{thm:all}
For every $n\ge3$ there is a fully indecomposable $A\in\Omega_n$ with exactly
$2n+3$ positive entries, each equal to $\frac14$, $\frac38$ or $\frac12$, such
that
\[
  \per A=\frac{15}{16}\cdot2^{1-n}.
\]
\end{theorem}

Chaining several dense blocks makes the failure exponential.

\begin{theorem}\label{thm:exp}
For every $n=4m$ with $m\ge3$ there is a fully indecomposable $A\in\Omega_n$
with $\frac{11}{4}n\le3n-3$ positive entries such that
\[
  \per A<(0.978)^{n}\cdot2^{1-n}.
\]
\end{theorem}

None of these matrices is nearly decomposable, in agreement with Lavi's
theorem. So in Foregger's inequality the hypothesis ``nearly decomposable''
cannot be weakened to ``fully indecomposable with at most $3n-3$ positive
entries''.

So $2n+3$ positive entries are too many. This is the exact threshold: a fully
indecomposable matrix has at least $2n$ positive entries, and with at most two
more the inequality survives.

\begin{proposition}\label{prop:threshold}
Every fully indecomposable $A\in\Omega_n$ with at most $2n+2$ positive entries
satisfies $\per A\ge2^{1-n}$.
\end{proposition}

The equality case for at most $2n+1$ positive entries is determined
in~\cite{FirstGap}. It follows from Proposition~\ref{prop:threshold} that the
order $6$ in Theorem~\ref{thm:six} cannot be lowered (Remark~\ref{rem:small}).

\section{The construction}

Fix integers $m\ge1$ and $k\ge0$ and a real number $0<q<1$, and put
\[
  a=\frac{1+q}{4},\qquad c=\frac{1-q}{2},\qquad
  B=\begin{pmatrix}a&a&c\\a&a&c\\c&c&0\end{pmatrix}.
\]
The row and column sums of $B$ are $1,1,1-q$. The matrix $N=N(m,k,q)$ has
order $n=m(k+3)$. Its rows and columns are divided into $m$ consecutive groups
of size $k+3$. In each group the first three rows and columns carry a copy of
$B$, called a \emph{block}, and the remaining $k$ rows and columns are called
\emph{ear} rows and columns. Besides the blocks, the positive entries of $N$
are the following, for each group:
\begin{itemize}
\item $1-q$ in position ($i$th ear row, $i$th ear column), for $1\le i\le k$;
\item $q$ in position ($i$th ear row, $(i+1)$st ear column), for $1\le i<k$;
\item $q$ in position (third block row, first ear column);
\item $q$ in position ($k$th ear row, third column of the block of the next
group), where the group after the last is the first.
\end{itemize}
If $k=0$, the last two items are replaced by one entry $q$ in position (third
block row, third column of the block of the next group). So each group has a
path, its \emph{ear}, that leaves the third row of its block, runs through its
ear columns and rows, and ends in the third column of the next block. The matrix
of Theorem~\ref{thm:six} is $N(1,3,\half)$.

\begin{lemma}\label{lem:necklace}
The matrix $N=N(m,k,q)$ is doubly stochastic and fully indecomposable, it is
not nearly decomposable, it has $m(2k+9)=2n+3m$ positive entries, and
\begin{equation}\label{eq:per}
  \per N=\Bigl(\frac{q^{k+1}(1+q)^2}{8}\Bigr)^{m}
        +\Bigl(\frac{(1-q)^{k+2}(1+q)}{4}\Bigr)^{m}.
\end{equation}
\end{lemma}

\begin{proof}
The first two rows of a block sum to $2a+c=1$, the third to $2c+q=1$, and an
ear row to $(1-q)+q=1$. The columns are the same, the third column of a block
receiving its $q$ from the previous ear. Each block has $8$ positive entries
and each ear $2k+1$.

For the permanent, consider a permutation all of whose entries in $N$ are
positive. Look at one ear. Its first column is covered either by the first ear
row or by the third block row. In the first case the second ear column must be
covered by the second ear row, and so on: all ear rows use their entries
$1-q$, and we call the ear \emph{idle}. In the second case the first ear row
must use its entry $q$, and so on: the ear uses all its $k+1$ entries $q$,
including the one in the third column of the next block, and we call it
\emph{active}. If $k=0$, the ear is active when its single entry $q$ is used and
idle otherwise.

If the ear leaving a block is active, the third row of the block is used up,
so the first two columns of the block are covered by its first two rows, and
its third column must be covered by the ear arriving at the block. That ear is
therefore active too. Going around the groups, either all ears are active or
all are idle. If all are active, each group contributes $q^{k+1}$ times the
permanent of the leading $2\times2$ part of $B$, which is $2a^2=(1+q)^2/8$. If
all are idle, each group contributes $(1-q)^{k}$ times
$\per B=4ac^2=(1+q)(1-q)^2/4$. This proves~\eqref{eq:per}.

A $(0,1)$-matrix of order $n\ge2$ is fully indecomposable if and only if its
bipartite graph is connected and every positive entry lies on a permutation
all of whose entries are positive \cite[Section~4.2]{BrualdiRyser}. The graph
of $N$ is connected. The permutations with all ears active cover every entry
$q$ and the four entries $a$ of each block, and those with all ears idle cover
every entry $1-q$ and every positive entry of each block. So $N$ is fully
indecomposable.

Finally, as T.~H. Foregger pointed out to the author, $N$ is not nearly
decomposable because each block contains a positive $2\times2$ submatrix: a
nearly decomposable matrix of order at least $3$ has none~\cite[Lemma~2.3]{Lavi}.
\end{proof}

\begin{proof}[Proof of Theorems~\ref{thm:six} and~\ref{thm:all}]
Take $m=1$, $q=\half$ and $k=n-3\ge0$. Then $a=\frac38$, $c=\frac14$, and the
number of positive entries is $2n+3$. By~\eqref{eq:per},
\[
  \per N=\frac{9}{32}\cdot2^{-(k+1)}+\frac{3}{8}\cdot2^{-(k+2)}
        =\frac{15}{64}\cdot2^{-k}=\frac{15}{16}\cdot2^{1-n}.\qedhere
\]
\end{proof}

\begin{proof}[Proof of Theorem~\ref{thm:exp}]
Take $k=1$ and $q=\frac{6}{13}$, so that $n=4m$ and $N$ has $11m$ positive
entries, which is at most $3n-3$ when $m\ge3$. By~\eqref{eq:per},
\[
  2^{n-1}\per N=\frac{\alpha^m+\beta^m}{2},\qquad
  \alpha=2q^2(1+q)^2=\frac{25992}{28561},\quad
  \beta=4(1-q)^3(1+q)=\frac{26068}{28561}.
\]
Since $\alpha<\beta$ and $\beta^{1/4}<0.978$, we get
$2^{n-1}\per N<\beta^{m}<(0.978)^{n}$.
\end{proof}

\section{The threshold}

A matrix of order $m\ge2$ is fully indecomposable if and only if every set of
$s$ rows, $0<s<m$, has positive entries in at least $s+1$ columns.

\begin{proof}[Proof of Proposition~\ref{prop:threshold}]
We use induction on $n$. For $n=2$ the permanent is $t^2+(1-t)^2\ge\half$. Let
$n\ge3$, and let $B$ minimize the permanent among the doubly stochastic matrices
that vanish wherever $A$ does. By Foregger's first
lemma~\cite[Lemma~1]{Foregger} (see~\cite[Lemma~2.1]{Lavi}), $B$ is fully
indecomposable and $\per B(i|j)=\per B$ whenever $b_{ij}>0$, where $B(i|j)$ is
$B$ with row $i$ and column $j$ deleted. Every row of $B$ has at least two
positive entries and $2n+2<3n$, so some row $i$ has exactly two, in columns $j$
and $k$. Delete row $i$ and replace columns $j$ and $k$ by their sum. The result
$C$ is doubly stochastic of order $n-1$ and has at most $2n=2(n-1)+2$ positive
entries. It is fully indecomposable: a set of rows violating the condition above
in $C$ violates it in $B$, after adding row $i$ if the set meets both columns
$j$ and $k$. The permanent is linear in the new column, so
$\per C=\per B(i|j)+\per B(i|k)=2\per B$. By induction,
$\per A\ge\per B=\half\per C\ge2^{1-n}$.
\end{proof}

\section{Remarks}

\begin{remark}[Why it works]
The block $B$ is close to the $3\times3$ matrix with all entries $\frac13$,
whose permanent is below $2^{1-3}$. A dense block costs three entries more
than the budget allows, and an ear of $k$ rows earns back $k$ of them, because
each of its rows has only two entries where the budget allows three. The ear
is cheap in permanent as well: it multiplies the permanent by about $2^{-k}$,
exactly as a path does in the extremal matrix $\half(I+C_n)$.
\end{remark}

\begin{remark}[Constants]
The value $q=\half$ was chosen for simple entries. For $n=6$ the best choice
in~\eqref{eq:per} is $q\approx0.4723$, which gives
$2^{5}\per N\approx0.9134$. In Theorem~\ref{thm:exp} the rate $0.978$ is not
optimal either. On the other hand it cannot be improved beyond $\frac89$: by
Gurvits's inequality \cite{Gurvits,LaurentSchrijver}, a doubly stochastic
matrix with $\Delta_j$ positive entries in column $j$ has
$\per A\ge\prod_j g(\Delta_j)$, where
$g(d)=\bigl(\frac{d-1}{d}\bigr)^{d-1}$. The numbers $-\log g(d)$ form an
increasing concave sequence, so if the $\Delta_j$ add up to less than $3n$
then $\per A\ge g(3)^n=(\frac49)^n$, that is,
$2^{n-1}\per A\ge\half(\frac89)^n$.
\end{remark}

\begin{remark}[Small orders]\label{rem:small}
Theorem~\ref{thm:six} has the smallest possible order among counterexamples
within the budget $3n-3$. A doubly stochastic matrix that is not fully
indecomposable is, up to permutations, a direct sum of smaller matrices, each of
order $1$ or fully indecomposable, and for $n\le5$ van der Waerden's bound on
the summands shows that it satisfies~\eqref{eq:foregger}. If $A$ is fully
indecomposable and $n\le5$, the budget $3n-3$ is at most $2n+2$, and
Proposition~\ref{prop:threshold} applies.
\end{remark}

\subsection*{Verification}
The ancillary script \texttt{verify.py} checks in exact arithmetic
Theorem~\ref{thm:six}, Theorem~\ref{thm:all} for $3\le n\le11$,
Theorem~\ref{thm:exp} for $m=3,4,5$, and Lemma~\ref{lem:necklace} on $40$
instances. It also checks the constants in the remark on constants.

\subsection*{Acknowledgment}
I thank T.~H. Foregger for comments on the construction.

\begin{thebibliography}{9}

\bibitem{BrualdiRyser}
R.~A. Brualdi and H.~J. Ryser, \emph{Combinatorial Matrix Theory}, Cambridge
University Press, 1991.

\bibitem{Foregger}
T.~H. Foregger, On the minimum value of the permanent of a nearly decomposable
doubly stochastic matrix, \emph{Linear Algebra Appl.} \textbf{32} (1980),
75--85.

\bibitem{Gurvits}
L.~Gurvits, Van der Waerden/Schrijver--Valiant like conjectures and stable
(aka hyperbolic) homogeneous polynomials: one theorem for all,
\emph{Electron. J. Combin.} \textbf{15} (2008), \#R66.

\bibitem{LaurentSchrijver}
M.~Laurent and A.~Schrijver, On Leonid Gurvits's proof for permanents,
\emph{Amer. Math. Monthly} \textbf{117} (2010), 903--911.

\bibitem{Lavi}
Y.~Lavi, Foregger's conjecture on minimum permanents of nearly decomposable
doubly stochastic matrices, preprint, 2026,
\href{https://doi.org/10.5281/zenodo.24208987}{doi:10.5281/zenodo.24208987}.

\bibitem{FirstGap}
T.~Maciejewski, One entry beyond the cycle: a sharp gap in Foregger's
inequality, preprint, 2026.

\bibitem{Minc}
H.~Minc, Nearly decomposable matrices, \emph{Linear Algebra Appl.} \textbf{5}
(1972), 181--187.

\end{thebibliography}

\end{document}
